% JobForge Resume Template -- XeLaTeX + CJK Chinese Support
% Based on LLM-Resume-Template & billryan/resume
\documentclass[11pt,a4paper]{article}

\usepackage[top=1.2cm,bottom=1.2cm,left=1.5cm,right=1.5cm]{geometry}
\usepackage{xeCJK}
\usepackage{fontspec}
\usepackage{titlesec}
\usepackage{enumitem}
\usepackage{hyperref}
\usepackage{xcolor}
\usepackage{tabularx}
\usepackage{fancyhdr}

\setCJKmainfont{STSong}
\setmainfont{Times New Roman}

\definecolor{primary}{RGB}{0, 82, 155}
\definecolor{darkgray}{RGB}{80, 80, 80}
\definecolor{lightgray}{RGB}{150, 150, 150}

\hypersetup{
    colorlinks=true,
    linkcolor=primary,
    urlcolor=primary,
}

\pagestyle{fancy}
\fancyhf{}
\renewcommand{\headrulewidth}{0pt}
\renewcommand{\footrulewidth}{0pt}

\titleformat{\section}
  {\Large\bfseries\color{primary}}
  {}
  {0em}
  {}
  [\titlerule]

\titlespacing*{\section}{0pt}{8pt}{4pt}

\setlist[itemize]{leftmargin=1.5em, itemsep=1pt, parsep=0pt, topsep=2pt}

\begin{document}

% ===== Header =====
\begin{center}
  {\Huge\bfseries （待补充）}\\[6pt]
  \textcolor{darkgray}{
    bcefghj@163.com \quad
    （待补充） \quad
    \href{https://github.com/bcefghj}{GitHub} \quad
    武汉
  }
\end{center}

\vspace{-4pt}

% ===== Summary =====
\section{个人总结}
专注于大模型应用与 Agent 架构研究，拥有从源码层面深度理解 Claude Code Agent 系统的独特优势。通过多个开源项目积累了 LLM 训练全流程、Agent 设计模式（ReAct/Tool Calling）、RAG 系统的工程落地经验，具备快速适配 LLM/Agent/RAG 任一方向的扎实基础。

% ===== Education =====

\section{教育背景}

\noindent\textbf{（待补充）} \hfill （待补充） \\
（待补充） · （待补充） 




% ===== Skills =====
\section{专业技能}
\begin{itemize}
  \item \textbf{编程语言}: Python, JavaScript, TypeScript, Rust
  \item \textbf{框架/工具}: PyTorch, HuggingFace Transformers, LangChain, LangGraph, Playwright
  \item \textbf{AI 知识}: LLM训练全流程 (Tokenizer/Pretrain/SFT/RLHF/DPO); Agent架构设计 (ReAct/Tool Calling/Memory/Planning/MCP协议); RAG系统 (向量检索/文档切分/Embedding/重排序); Claude Code 架构深度理解; Prompt Engineering; Transformer架构
  \item \textbf{开发工具}: Git/GitHub, Docker, Linux, LaTeX
\end{itemize}

% ===== Projects =====
\section{项目经历}

\noindent\textbf{miniClaudeCode - Agent架构最小复现} \hfill \textcolor{lightgray}{2026.03 - 至今}\\
\textit{从50万行 Claude Code 源码蒸馏为 ~1000 行 Python 的核心 Agent 架构最小复现}\\
技术栈: Python, Agent Framework, Tool Calling, ReAct, Streaming
\begin{itemize}

  \item 作为唯一开发者，从 50 万行源码中提炼 Claude Code 核心架构，实现 Tool Calling / Permission / Streaming 三大核心模块，为金山 Agent 研发提供可直接参考的架构范式

  \item 采用 ReAct 范式构建 Agent Loop，完整实现 Function Calling 解析与执行流程，验证了多轮推理+工具调用在复杂任务中的有效性

  \item 代码可直接作为企业级 Agent 开发起点，已获 40+ Stars，体现对 Agent 架构的深度理解与工程能力

\end{itemize}


\noindent\textbf{AI Agent 面试全攻略} \hfill \textcolor{lightgray}{2026.03 - 至今}\\
\textit{覆盖 200+ 面试题、企业级项目、简历模板的一站式 AI Agent 面试指南}\\
技术栈: Python, LLM, Agent, RAG, Java, Go
\begin{itemize}

  \item 针对金山 RAG/Agent 方向需求，系统整理了 Tool Calling、Memory 机制、Planning 策略等核心面试要点，帮助快速掌握 Agent 工程落地的关键技术栈

  \item 实现 Python/Java/Go 多语言企业级项目，涵盖 RAG 检索增强、Agent 任务规划等实战场景，与金山 LLM 应用算法研发方向高度匹配

  \item 采用 STAR 法则设计面试稿模板，体现对技术表达与问题拆解能力的刻意练习

\end{itemize}


\noindent\textbf{Learn MiniMind - LLM训练全流程} \hfill \textcolor{lightgray}{2026.03 - 至今}\\
\textit{从 Tokenizer 到 DPO，系统掌握 LLM 训练全流程的 22 节实战课程}\\
技术栈: PyTorch, Transformer, RLHF, DPO, HuggingFace
\begin{itemize}

  \item 完整实现预训练(Pretrain)、监督微调(SFT)、人类反馈强化学习(RLHF)、直接偏好优化(DPO)全流程，与金山 LLM 算法研发岗位的核心技术要求直接对应

  \item 深入理解 Transformer 架构与 PyTorch 实现，能够从底层机制解释大模型行为，为算法优化提供理论支撑

  \item 支持多模态扩展(learn-minimind-multimodal)，覆盖 CV/多模态 Agent 等前沿方向

\end{itemize}


\noindent\textbf{Claude Code 完全指南} \hfill \textcolor{lightgray}{2026.03 - 至今}\\
\textit{哆啦A梦风格漫画 + 详细源码解析的 Claude Code 完整指南（12篇120节课）}\\
技术栈: JavaScript, Agent架构, 源码分析, TypeScript
\begin{itemize}

  \item 独立完成 50 万行 Claude Code 源码的系统性架构分析，涵盖 Agent Loop、Permission System、Tool System 等核心模块，深度远超社区同类资源

  \item 以零基础友好方式讲解复杂架构，体现将复杂技术简单化、产品化的能力，有助于金山技术文档与用户指引的优化

  \item 获得 250+ GitHub Stars，成为该领域最受欢迎的中文教程之一，验证了技术研究与内容创作的双重能力

\end{itemize}


\noindent\textbf{企业级 RAG 系统实现（规划中）} \hfill \textcolor{lightgray}{规划中}\\
\textit{基于 LangChain/LangGraph 的生产级 RAG 系统，支持多路召回、重排序与 Agent 化扩展}\\
技术栈: LangChain, LangGraph, HuggingFace, RAG, 向量检索, Python
\begin{itemize}

  \item 规划实现 BM25 + 向量混合检索、Cross-Encoder 重排序、冷启动知识库构建等企业级 RAG 核心模块，直接对接金山 RAG 方向算法落地需求

  \item 计划集成 Agent 能力，实现 '检索->判断->再检索' 的自适应 RAG 流程，体现 RAG + Agent 融合的前沿技术视野

  \item 将使用 LangChain/LangGraph + HuggingFace Embedding + Qdrant/FAISS，构建完整的向量检索与知识管理 pipeline

\end{itemize}


\noindent\textbf{Agent 任务规划框架（规划中）} \hfill \textcolor{lightgray}{规划中}\\
\textit{基于 ReAct/ReWOO/Plan-and-Execute 等范式的可观测 Agent 任务规划框架}\\
技术栈: Python, Agent, Planning, LangGraph, PyTorch
\begin{itemize}

  \item 规划实现多策略 Planning（LLM-Based / Hierarchical / Tree-of-Thought），对比验证不同规划范式在复杂任务中的效果差异，为金山 Agent 算法选型提供参考

  \item 设计完整的 Agent 可观测性方案（Trace/Memory/Reflection），支持任务回放与失败诊断，提升 Agent 系统调试效率

  \item 将作为 miniClaudeCode 的上层扩展，实现从最小复现到生产级框架的演进

\end{itemize}



% ===== Keywords =====

\section{关键词}
\textcolor{lightgray}{\small LLM · Agent · RAG · 大模型 · 深度学习 · 算法研发 · PyTorch · HuggingFace · LangChain · LangGraph · Tool Calling · ReAct · RLHF · DPO · Transformer · 向量检索 · Embedding · MCP协议 · Prompt Engineering · 实习生}


\end{document}